\chapter*{\textbf{ПРИЛОЖЕНИЕ C}}
\addcontentsline{toc}{chapter}{\textbf{ПРИЛОЖЕНИЕ C}}


\begin{lstlisting}[
    caption={Сгенерированный LLVM IR код для тестовой программы №1},
    label={lst:while-ll},
    language={LLVM},
    basicstyle=\footnotesize\ttfamily,
    columns=fullflexible,
    keepspaces=true,
    breaklines=true
]

; ModuleID = 'generated'
source_filename = "generated"
target datalayout = "e-m:o-p270:32:32-p271:32:32-p272:64:64-i64:64-i128:128-n32:64-S128-Fn32"
target triple = "arm64-apple-darwin23.4.0"

@0 = private unnamed_addr constant [4 x i8] c"abc\00", align 1
@1 = private unnamed_addr constant [4 x i8] c"%d\0A\00", align 1
@2 = private unnamed_addr constant [4 x i8] c"%d\0A\00", align 1
@3 = private unnamed_addr constant [4 x i8] c"%s\0A\00", align 1
@4 = private unnamed_addr constant [4 x i8] c"%d\0A\00", align 1
@5 = private unnamed_addr constant [5 x i8] c"abcd\00", align 1
@6 = private unnamed_addr constant [4 x i8] c"%s\0A\00", align 1
@7 = private unnamed_addr constant [4 x i8] c"%d\0A\00", align 1
@8 = private unnamed_addr constant [4 x i8] c"%d\0A\00", align 1

define i32 @main() {
entry:
  %retval = alloca i32, align 4
  %a = alloca i32, align 4
  store i32 9, ptr %a, align 4
  %loaded_from_var = load i32, ptr %a, align 4
  store i32 %loaded_from_var, ptr %a, align 4
  %b = alloca i32, align 4
  store i32 34, ptr %b, align 4
  %loaded_from_var1 = load i32, ptr %b, align 4
  store i32 %loaded_from_var1, ptr %b, align 4
  %loaded_int = load i32, ptr %a, align 4
  call void (ptr, ...) @printf(ptr @1, i32 %loaded_int)
  %loaded_int2 = load i32, ptr %b, align 4
  call void (ptr, ...) @printf(ptr @2, i32 %loaded_int2)
  call void (ptr, ...) @printf(ptr @3, ptr @0)
  call void (ptr, ...) @printf(ptr @4, i32 2)
  call void (ptr, ...) @printf(ptr @6, ptr @5)
  %loaded_int3 = load i32, ptr %a, align 4
  call void (ptr, ...) @printf(ptr @7, i32 %loaded_int3)
  %loaded_int4 = load i32, ptr %b, align 4
  call void (ptr, ...) @printf(ptr @8, i32 %loaded_int4)
  %ret = load i32, ptr %retval, align 4
  ret i32 %ret
}

declare void @printf(ptr, ...)

\end{lstlisting}



\begin{lstlisting}[
    caption={Сгенерированный LLVM IR код для тестовой программы №2},
    label={lst:while-ll},
    language={LLVM},
    basicstyle=\footnotesize\ttfamily,
    columns=fullflexible,
    keepspaces=true,
    breaklines=true
]

; ModuleID = 'generated'
source_filename = "generated"
target datalayout = "e-m:o-p270:32:32-p271:32:32-p272:64:64-i64:64-i128:128-n32:64-S128-Fn32"
target triple = "arm64-apple-darwin23.4.0"

@0 = private unnamed_addr constant [4 x i8] c"%d\0A\00", align 1
@1 = private unnamed_addr constant [4 x i8] c"%d\0A\00", align 1
@2 = private unnamed_addr constant [4 x i8] c"%d\0A\00", align 1
@3 = private unnamed_addr constant [4 x i8] c"%d\0A\00", align 1
@4 = private unnamed_addr constant [4 x i8] c"%d\0A\00", align 1
@5 = private unnamed_addr constant [4 x i8] c"%d\0A\00", align 1

define i32 @main() {
entry:
  %retval = alloca i32, align 4
  %a = alloca i32, align 4
  store i32 6, ptr %a, align 4
  %loaded_from_var = load i32, ptr %a, align 4
  store i32 %loaded_from_var, ptr %a, align 4
  %b = alloca i32, align 4
  store i32 1981, ptr %b, align 4
  %loaded_from_var1 = load i32, ptr %b, align 4
  store i32 %loaded_from_var1, ptr %b, align 4
  %calltmp = call i32 @max(ptr %a, ptr %b)
  %tmp = alloca i32, align 4
  store i32 %calltmp, ptr %tmp, align 4
  %calltmp2 = call i32 @max(ptr %a, ptr %b)
  %ret = load i32, ptr %retval, align 4
  ret i32 %ret
}

define i32 @max(i32 %a, i32 %b) {
entry:
  %a1 = alloca i32, align 4
  store i32 %a, ptr %a1, align 4
  %b2 = alloca i32, align 4
  store i32 %b, ptr %b2, align 4
  %retval = alloca i32, align 4
  %tmp = alloca i32, align 4
  store i32 0, ptr %tmp, align 4
  %loaded_rhs = load i32, ptr %tmp, align 4
  %loaded_lhs = load i32, ptr %a1, align 4
  %gttmp = icmp sgt i32 %loaded_lhs, %loaded_rhs
  %ifcond = icmp ne i1 %gttmp, false
  br i1 %ifcond, label %if.then, label %if.end

if.then:                                          ; preds = %entry
  %tmp3 = alloca i32, align 4
  store i32 1, ptr %tmp3, align 4
  %loaded_int = load i32, ptr %tmp3, align 4
  call void (ptr, ...) @printf(ptr @0, i32 %loaded_int)
  br label %if.end

if.end:                                           ; preds = %if.then, %entry
  %a4 = alloca i32, align 4
  store i32 3, ptr %a4, align 4
  %loaded_from_var = load i32, ptr %a4, align 4
  store i32 %loaded_from_var, ptr %a4, align 4
  %result = alloca i32, align 4
  %loaded_rhs5 = load i32, ptr %a4, align 4
  store i32 %loaded_rhs5, ptr %result, align 4
  %loaded_int6 = load i32, ptr %result, align 4
  call void (ptr, ...) @printf(ptr @1, i32 %loaded_int6)
  %loaded_int7 = load i32, ptr %a4, align 4
  call void (ptr, ...) @printf(ptr @2, i32 %loaded_int7)
  %ret_val = load i32, ptr %result, align 4
  ret i32 %ret_val

entry8:                                           ; No predecessors!
  %a9 = alloca i32, align 4
  store i32 %a, ptr %a9, align 4
  %b10 = alloca i32, align 4
  store i32 %b, ptr %b10, align 4
  %retval11 = alloca i32, align 4
  %loaded_lhs12 = load i32, ptr %a9, align 4
  %gttmp13 = icmp sgt i32 %loaded_lhs12, 0
  %ifcond14 = icmp ne i1 %gttmp13, false
  br i1 %ifcond14, label %if.then15, label %if.end16

if.then15:                                        ; preds = %entry8
  call void (ptr, ...) @printf(ptr @3, i32 1)
  br label %if.end16

if.end16:                                         ; preds = %if.then15, %entry8
  %a17 = alloca i32, align 4
  store i32 3, ptr %a17, align 4
  %loaded_from_var18 = load i32, ptr %a17, align 4
  store i32 %loaded_from_var18, ptr %a17, align 4
  %result19 = alloca i32, align 4
  %loaded_rhs20 = load i32, ptr %a17, align 4
  store i32 %loaded_rhs20, ptr %result19, align 4
  %loaded_int21 = load i32, ptr %result19, align 4
  call void (ptr, ...) @printf(ptr @4, i32 %loaded_int21)
  %loaded_int22 = load i32, ptr %a17, align 4
  call void (ptr, ...) @printf(ptr @5, i32 %loaded_int22)
  %ret_val23 = load i32, ptr %result19, align 4
  ret i32 %ret_val23
}

declare void @printf(ptr, ...)

\end{lstlisting}


\begin{lstlisting}[
    caption={Сгенерированный LLVM IR код для тестовой программы №3},
    label={lst:while-ll},
    language={LLVM},
    basicstyle=\footnotesize\ttfamily,
    columns=fullflexible,
    keepspaces=true,
    breaklines=true
]

; ModuleID = 'generated'
source_filename = "generated"
target datalayout = "e-m:o-p270:32:32-p271:32:32-p272:64:64-i64:64-i128:128-n32:64-S128-Fn32"
target triple = "arm64-apple-darwin23.4.0"

@0 = private unnamed_addr constant [4 x i8] c"%d\0A\00", align 1
@1 = private unnamed_addr constant [4 x i8] c"%d\0A\00", align 1

define i32 @main() {
entry:
  %retval = alloca i32, align 4
  %i = alloca i32, align 4
  store i32 0, ptr %i, align 4
  %loaded_from_var = load i32, ptr %i, align 4
  store i32 %loaded_from_var, ptr %i, align 4
  br label %while.cond

while.cond:                                       ; preds = %while.body, %entry
  %loaded_lhs = load i32, ptr %i, align 4
  %letmp = icmp sle i32 %loaded_lhs, 10
  %whilecond = icmp ne i1 %letmp, false
  br i1 %whilecond, label %while.body, label %while.end

while.body:                                       ; preds = %while.cond
  %loaded_int = load i32, ptr %i, align 4
  call void (ptr, ...) @printf(ptr @0, i32 %loaded_int)
  %loadtmp = load i32, ptr %i, align 4
  %inctmp = add i32 %loadtmp, 1
  store i32 %inctmp, ptr %i, align 4
  %loadtmp1 = load i32, ptr %i, align 4
  %inctmp2 = add i32 %loadtmp1, 1
  store i32 %inctmp2, ptr %i, align 4
  br label %while.cond

while.end:                                        ; preds = %while.cond
  %j = alloca i32, align 4
  store i32 0, ptr %j, align 4
  %loaded_from_var3 = load i32, ptr %j, align 4
  store i32 %loaded_from_var3, ptr %j, align 4
  br label %while.cond4

while.cond4:                                      ; preds = %while.body5, %while.end
  %loaded_lhs7 = load i32, ptr %j, align 4
  %letmp8 = icmp sle i32 %loaded_lhs7, 10
  %whilecond9 = icmp ne i1 %letmp8, false
  br i1 %whilecond9, label %while.body5, label %while.end6

while.body5:                                      ; preds = %while.cond4
  %loaded_int10 = load i32, ptr %j, align 4
  call void (ptr, ...) @printf(ptr @1, i32 %loaded_int10)
  %loadtmp11 = load i32, ptr %j, align 4
  %inctmp12 = add i32 %loadtmp11, 1
  store i32 %inctmp12, ptr %j, align 4
  br label %while.cond4

while.end6:                                       ; preds = %while.cond4
  %k = alloca i32, align 4
  store i32 0, ptr %k, align 4
  %loaded_from_var13 = load i32, ptr %k, align 4
  store i32 %loaded_from_var13, ptr %k, align 4
  %ret = load i32, ptr %retval, align 4
  ret i32 %ret
}

declare void @printf(ptr, ...)

\end{lstlisting}



\begin{lstlisting}[
    caption={Сгенерированный LLVM IR код для тестовой программы №4},
    label={lst:while-ll},
    language={LLVM},
    basicstyle=\footnotesize\ttfamily,
    columns=fullflexible,
    keepspaces=true,
    breaklines=true
]

; ModuleID = 'generated'
source_filename = "generated"
target datalayout = "e-m:o-p270:32:32-p271:32:32-p272:64:64-i64:64-i128:128-n32:64-S128-Fn32"
target triple = "arm64-apple-darwin23.4.0"

@0 = private unnamed_addr constant [4 x i8] c"%d\0A\00", align 1

define i32 @main() {
entry:
  %retval = alloca i32, align 4
  %x = alloca i32, align 4
  store i32 1, ptr %x, align 4
  %loaded_from_var = load i32, ptr %x, align 4
  store i32 %loaded_from_var, ptr %x, align 4
  %loadtmp = load i32, ptr %x, align 4
  %inctmp = add i32 %loadtmp, 1
  store i32 %inctmp, ptr %x, align 4
  %loadtmp1 = load i32, ptr %x, align 4
  %inctmp2 = add i32 %loadtmp1, 1
  store i32 %inctmp2, ptr %x, align 4
  %loadtmp3 = load i32, ptr %x, align 4
  %dectmp = sub i32 %loadtmp3, 1
  store i32 %dectmp, ptr %x, align 4
  %loaded_int = load i32, ptr %x, align 4
  call void (ptr, ...) @printf(ptr @0, i32 %loaded_int)
  %ret = load i32, ptr %retval, align 4
  ret i32 %ret
}

declare void @printf(ptr, ...)


\end{lstlisting}



\begin{lstlisting}[
    caption={Сгенерированный LLVM IR код для тестовой программы №5},
    label={lst:while-ll},
    language={LLVM},
    basicstyle=\footnotesize\ttfamily,
    columns=fullflexible,
    keepspaces=true,
    breaklines=true
]

; ModuleID = 'generated'
source_filename = "generated"
target datalayout = "e-m:o-p270:32:32-p271:32:32-p272:64:64-i64:64-i128:128-n32:64-S128-Fn32"
target triple = "arm64-apple-darwin23.4.0"

@0 = private unnamed_addr constant [4 x i8] c"%d\0A\00", align 1
@1 = private unnamed_addr constant [4 x i8] c"%d\0A\00", align 1
@2 = private unnamed_addr constant [4 x i8] c"%d\0A\00", align 1

define i32 @main() {
entry:
  %retval = alloca i32, align 4
  %n = alloca i32, align 4
  store i32 5, ptr %n, align 4
  %loaded_from_var = load i32, ptr %n, align 4
  store i32 %loaded_from_var, ptr %n, align 4
  br label %while.cond

while.cond:                                       ; preds = %while.body, %entry
  %loaded_lhs = load i32, ptr %n, align 4
  %netmp = icmp ne i32 %loaded_lhs, 0
  %whilecond = icmp ne i1 %netmp, false
  br i1 %whilecond, label %while.body, label %while.end

while.body:                                       ; preds = %while.cond
  %loaded_int = load i32, ptr %n, align 4
  call void (ptr, ...) @printf(ptr @0, i32 %loaded_int)
  %loadtmp = load i32, ptr %n, align 4
  %dectmp = sub i32 %loadtmp, 1
  store i32 %dectmp, ptr %n, align 4
  br label %while.cond

while.end:                                        ; preds = %while.cond
  %calltmp = call i32 @f(i32 24, i32 42)
  ret i32 0
}

define i32 @f(i32 %s, i32 %y) {
entry:
  %s1 = alloca i32, align 4
  store i32 %s, ptr %s1, align 4
  %y2 = alloca i32, align 4
  store i32 %y, ptr %y2, align 4
  %retval = alloca i32, align 4
  %loaded_int = load i32, ptr %s1, align 4
  call void (ptr, ...) @printf(ptr @1, i32 %loaded_int)
  %loaded_int3 = load i32, ptr %y2, align 4
  call void (ptr, ...) @printf(ptr @2, i32 %loaded_int3)
  ret i32 0
}

declare void @printf(ptr, ...)


\end{lstlisting}


\begin{lstlisting}[
    caption={Сгенерированный LLVM IR код для тестовой программы №6},
    label={lst:while-ll},
    language={LLVM},
    basicstyle=\footnotesize\ttfamily,
    columns=fullflexible,
    keepspaces=true,
    breaklines=true
]

; ModuleID = 'generated'
source_filename = "generated"
target datalayout = "e-m:o-p270:32:32-p271:32:32-p272:64:64-i64:64-i128:128-n32:64-S128-Fn32"
target triple = "arm64-apple-darwin23.4.0"

@0 = private unnamed_addr constant [4 x i8] c"Hi!\00", align 1
@1 = private unnamed_addr constant [4 x i8] c"%s\0A\00", align 1
@2 = private unnamed_addr constant [4 x i8] c"%d\0A\00", align 1
@3 = private unnamed_addr constant [4 x i8] c"%d\0A\00", align 1
@4 = private unnamed_addr constant [4 x i8] c"%d\0A\00", align 1

define i32 @main() {
entry:
  %retval = alloca i32, align 4
  %calltmp = call i32 @Print1()
  %calltmp1 = call i32 @Print2(i32 31415)
  %calltmp2 = call i32 @Print3(i32 777, i32 34)
  %ret = load i32, ptr %retval, align 4
  ret i32 %ret
}

define i32 @Print1() {
entry:
  %retval = alloca i32, align 4
  call void (ptr, ...) @printf(ptr @1, ptr @0)
  %ret = load i32, ptr %retval, align 4
  ret i32 %ret
}

define i32 @Print2(i32 %x) {
entry:
  %x1 = alloca i32, align 4
  store i32 %x, ptr %x1, align 4
  %retval = alloca i32, align 4
  %loaded_int = load i32, ptr %x1, align 4
  call void (ptr, ...) @printf(ptr @2, i32 %loaded_int)
  %ret = load i32, ptr %retval, align 4
  ret i32 %ret
}

define i32 @Print3(i32 %x, i32 %y) {
entry:
  %x1 = alloca i32, align 4
  store i32 %x, ptr %x1, align 4
  %y2 = alloca i32, align 4
  store i32 %y, ptr %y2, align 4
  %retval = alloca i32, align 4
  %loaded_int = load i32, ptr %x1, align 4
  call void (ptr, ...) @printf(ptr @3, i32 %loaded_int)
  %loaded_int3 = load i32, ptr %y2, align 4
  call void (ptr, ...) @printf(ptr @4, i32 %loaded_int3)
  %ret = load i32, ptr %retval, align 4
  ret i32 %ret
}

declare void @printf(ptr, ...)


\end{lstlisting}



\begin{lstlisting}[
    caption={Сгенерированный LLVM IR код для тестовой программы №7},
    label={lst:while-ll},
    language={LLVM},
    basicstyle=\footnotesize\ttfamily,
    columns=fullflexible,
    keepspaces=true,
    breaklines=true
]

; ModuleID = 'generated'
source_filename = "generated"
target datalayout = "e-m:o-p270:32:32-p271:32:32-p272:64:64-i64:64-i128:128-n32:64-S128-Fn32"
target triple = "arm64-apple-darwin23.4.0"

@0 = private unnamed_addr constant [4 x i8] c"%d\0A\00", align 1
@1 = private unnamed_addr constant [4 x i8] c"%d\0A\00", align 1
@2 = private unnamed_addr constant [4 x i8] c"%d\0A\00", align 1

define i32 @max(i32 %a, i32 %b) {
entry:
  %a1 = alloca i32, align 4
  store i32 %a, ptr %a1, align 4
  %b2 = alloca i32, align 4
  store i32 %b, ptr %b2, align 4
  %retval = alloca i32, align 4
  %loaded_cond = load i32, ptr %b2, align 4
  %cond_bool = icmp ne i32 %loaded_cond, 0
  br i1 %cond_bool, label %ternary.true, label %ternary.false

ternary.true:                                     ; preds = %entry
  br label %ternary.merge

ternary.false:                                    ; preds = %entry
  br label %ternary.merge

ternary.merge:                                    ; preds = %ternary.false, %ternary.true
  %ternarytmp = phi ptr [ %a1, %ternary.true ], [ %b2, %ternary.false ]
  %loaded_lhs = load i32, ptr %a1, align 4
  %loaded_rhs = load i32, ptr %ternarytmp, align 4
  %gttmp = icmp sgt i32 %loaded_lhs, %loaded_rhs
  %select_result = select i1 %gttmp, ptr %a1, ptr %b2
  %selected_loaded = load i32, ptr %select_result, align 4
  %result = alloca i32, align 4
  store i32 %selected_loaded, ptr %result, align 4
  %loaded_int = load i32, ptr %a1, align 4
  call void (ptr, ...) @printf(ptr @0, i32 %loaded_int)
  %loaded_int3 = load i32, ptr %b2, align 4
  call void (ptr, ...) @printf(ptr @1, i32 %loaded_int3)
  %loaded_int4 = load i32, ptr %result, align 4
  call void (ptr, ...) @printf(ptr @2, i32 %loaded_int4)
  %ret = load i32, ptr %retval, align 4
  ret i32 %ret
}

define i32 @main() {
entry:
  %retval = alloca i32, align 4
  %calltmp = call i32 @max(i32 10, i32 20)
  %ret = load i32, ptr %retval, align 4
  ret i32 %ret
}

declare void @printf(ptr, ...)


\end{lstlisting}



\begin{lstlisting}[
    caption={Сгенерированный LLVM IR код для тестовой программы №8},
    label={lst:while-ll},
    language={LLVM},
    basicstyle=\footnotesize\ttfamily,
    columns=fullflexible,
    keepspaces=true,
    breaklines=true
]

; ModuleID = 'generated'
source_filename = "generated"
target datalayout = "e-m:o-p270:32:32-p271:32:32-p272:64:64-i64:64-i128:128-n32:64-S128-Fn32"
target triple = "arm64-apple-darwin23.4.0"

@0 = private unnamed_addr constant [10 x i8] c"a + b > 0\00", align 1
@1 = private unnamed_addr constant [4 x i8] c"%s\0A\00", align 1
@2 = private unnamed_addr constant [11 x i8] c"a + b <= 0\00", align 1
@3 = private unnamed_addr constant [4 x i8] c"%s\0A\00", align 1
@4 = private unnamed_addr constant [10 x i8] c"a + b > 0\00", align 1
@5 = private unnamed_addr constant [4 x i8] c"%s\0A\00", align 1
@6 = private unnamed_addr constant [11 x i8] c"a + b <= 0\00", align 1
@7 = private unnamed_addr constant [4 x i8] c"%s\0A\00", align 1

define i32 @main() {
entry:
  %retval = alloca i32, align 4
  %calltmp = call i32 @someFunc(i32 1982, i32 3)
  %calltmp1 = call i32 @someFunc(i32 0, i32 0)
  %ret = load i32, ptr %retval, align 4
  ret i32 %ret
}

define i32 @someFunc(i32 %a, i32 %b) {
entry:
  %a1 = alloca i32, align 4
  store i32 %a, ptr %a1, align 4
  %b2 = alloca i32, align 4
  store i32 %b, ptr %b2, align 4
  %retval = alloca i32, align 4
  %loaded_lhs = load i32, ptr %a1, align 4
  %loaded_rhs = load i32, ptr %b2, align 4
  %addtmp = add i32 %loaded_lhs, %loaded_rhs
  %gttmp = icmp sgt i32 %addtmp, 0
  %ifcond = icmp ne i1 %gttmp, false
  br i1 %ifcond, label %if.then, label %if.else

if.then:                                          ; preds = %entry
  call void (ptr, ...) @printf(ptr @1, ptr @0)
  br label %if.end

if.end:                                           ; preds = %if.else, %if.then
  %ret = load i32, ptr %retval, align 4
  ret i32 %ret

if.else:                                          ; preds = %entry
  call void (ptr, ...) @printf(ptr @3, ptr @2)
  br label %if.end

entry3:                                           ; No predecessors!
  %a4 = alloca i32, align 4
  store i32 %a, ptr %a4, align 4
  %b5 = alloca i32, align 4
  store i32 %b, ptr %b5, align 4
  %retval6 = alloca i32, align 4
  %loaded_lhs7 = load i32, ptr %a4, align 4
  %loaded_rhs8 = load i32, ptr %b5, align 4
  %addtmp9 = add i32 %loaded_lhs7, %loaded_rhs8
  %gttmp10 = icmp sgt i32 %addtmp9, 0
  %ifcond11 = icmp ne i1 %gttmp10, false
  br i1 %ifcond11, label %if.then12, label %if.else14

if.then12:                                        ; preds = %entry3
  call void (ptr, ...) @printf(ptr @5, ptr @4)
  br label %if.end13

if.end13:                                         ; preds = %if.else14, %if.then12
  %ret15 = load i32, ptr %retval6, align 4
  ret i32 %ret15

if.else14:                                        ; preds = %entry3
  call void (ptr, ...) @printf(ptr @7, ptr @6)
  br label %if.end13
}

declare void @printf(ptr, ...)

\end{lstlisting}



\begin{lstlisting}[
    caption={Сгенерированный LLVM IR код для тестовой программы №9},
    label={lst:while-ll},
    language={LLVM},
    basicstyle=\footnotesize\ttfamily,
    columns=fullflexible,
    keepspaces=true,
    breaklines=true
]

; ModuleID = 'generated'
source_filename = "generated"
target datalayout = "e-m:o-p270:32:32-p271:32:32-p272:64:64-i64:64-i128:128-n32:64-S128-Fn32"
target triple = "arm64-apple-darwin23.4.0"

@0 = private unnamed_addr constant [4 x i8] c"%d\0A\00", align 1
@1 = private unnamed_addr constant [4 x i8] c"%d\0A\00", align 1
@2 = private unnamed_addr constant [4 x i8] c"%d\0A\00", align 1

define i32 @main() {
entry:
  %retval = alloca i32, align 4
  %x = alloca i32, align 4
  store i32 2, ptr %x, align 4
  %loaded_from_var = load i32, ptr %x, align 4
  store i32 %loaded_from_var, ptr %x, align 4
  %switch_cond = load i32, ptr %x, align 4
  switch i32 %switch_cond, label %switch.end [
    i32 1, label %case
    i32 2, label %case1
    i32 3, label %case2
  ]

switch.end:                                       ; preds = %case2, %case1, %case, %entry
  %ret = load i32, ptr %retval, align 4
  ret i32 %ret

case:                                             ; preds = %entry
  call void (ptr, ...) @printf(ptr @0, i32 100)
  br label %switch.end

case1:                                            ; preds = %entry
  call void (ptr, ...) @printf(ptr @1, i32 200)
  br label %switch.end

case2:                                            ; preds = %entry
  call void (ptr, ...) @printf(ptr @2, i32 300)
  br label %switch.end
}

declare void @printf(ptr, ...)


\end{lstlisting}


\begin{lstlisting}[
    caption={Сгенерированный LLVM IR код для тестовой программы №10},
    label={lst:while-ll},
    language={LLVM},
    basicstyle=\footnotesize\ttfamily,
    columns=fullflexible,
    keepspaces=true,
    breaklines=true
]

; ModuleID = 'generated'
source_filename = "generated"
target datalayout = "e-m:o-p270:32:32-p271:32:32-p272:64:64-i64:64-i128:128-n32:64-S128-Fn32"
target triple = "arm64-apple-darwin23.4.0"

@0 = private unnamed_addr constant [4 x i8] c"%d\0A\00", align 1
@1 = private unnamed_addr constant [4 x i8] c"%d\0A\00", align 1
@2 = private unnamed_addr constant [4 x i8] c"%d\0A\00", align 1
@3 = private unnamed_addr constant [4 x i8] c"%d\0A\00", align 1

define i32 @main() {
entry:
  %retval = alloca i32, align 4
  %x = alloca i32, align 4
  store i32 5, ptr %x, align 4
  %loaded_from_var = load i32, ptr %x, align 4
  store i32 %loaded_from_var, ptr %x, align 4
  %calltmp = call i32 @goTOfunc(i32 0)
  %x1 = alloca i32, align 4
  store i32 3, ptr %x1, align 4
  %loaded_from_var2 = load i32, ptr %x1, align 4
  store i32 %loaded_from_var2, ptr %x1, align 4
  %calltmp3 = call i32 @goTOfunc(i32 4)
  %ret = load i32, ptr %retval, align 4
  ret i32 %ret
}

define i32 @goTOfunc(i32 %y) {
entry:
  %y1 = alloca i32, align 4
  store i32 %y, ptr %y1, align 4
  %retval = alloca i32, align 4
  %loaded_lhs = load i32, ptr %y1, align 4
  %gttmp = icmp sgt i32 %loaded_lhs, 0
  %ifcond = icmp ne i1 %gttmp, false
  br i1 %ifcond, label %if.then, label %if.end

if.then:                                          ; preds = %entry
  br label %label1

if.end:                                           ; preds = %after_goto, %entry
  call void (ptr, ...) @printf(ptr @0, i32 0)
  br label %label1

label1:                                           ; preds = %if.end9, %if.then8, %if.end, %if.then
  call void (ptr, ...) @printf(ptr @1, i32 42)
  %ret = load i32, ptr %retval, align 4
  ret i32 %ret
  call void (ptr, ...) @printf(ptr @3, i32 42)
  %ret11 = load i32, ptr %retval4, align 4
  ret i32 %ret11

after_goto:                                       ; No predecessors!
  br label %if.end

entry2:                                           ; No predecessors!
  %y3 = alloca i32, align 4
  store i32 %y, ptr %y3, align 4
  %retval4 = alloca i32, align 4
  %loaded_lhs5 = load i32, ptr %y3, align 4
  %gttmp6 = icmp sgt i32 %loaded_lhs5, 0
  %ifcond7 = icmp ne i1 %gttmp6, false
  br i1 %ifcond7, label %if.then8, label %if.end9

if.then8:                                         ; preds = %entry2
  br label %label1

if.end9:                                          ; preds = %after_goto10, %entry2
  call void (ptr, ...) @printf(ptr @2, i32 0)
  br label %label1

after_goto10:                                     ; No predecessors!
  br label %if.end9
}

declare void @printf(ptr, ...)

\end{lstlisting}


\begin{lstlisting}[
    caption={Сгенерированный LLVM IR код для тестовой программы №11},
    label={lst:while-ll},
    language={LLVM},
    basicstyle=\footnotesize\ttfamily,
    columns=fullflexible,
    keepspaces=true,
    breaklines=true
]

; ModuleID = 'generated'
source_filename = "generated"
target datalayout = "e-m:o-p270:32:32-p271:32:32-p272:64:64-i64:64-i128:128-n32:64-S128-Fn32"
target triple = "arm64-apple-darwin23.4.0"

@0 = private unnamed_addr constant [4 x i8] c"%d\0A\00", align 1

define i32 @main() {
entry:
  %retval = alloca i32, align 4
  %max = alloca i32, align 4
  store i32 2, ptr %max, align 4
  %loaded_from_var = load i32, ptr %max, align 4
  store i32 %loaded_from_var, ptr %max, align 4
  %x = alloca i32, align 4
  store i32 10, ptr %x, align 4
  %loaded_from_var1 = load i32, ptr %x, align 4
  store i32 %loaded_from_var1, ptr %x, align 4
  %y = alloca i32, align 4
  store i32 20, ptr %y, align 4
  %loaded_from_var2 = load i32, ptr %y, align 4
  store i32 %loaded_from_var2, ptr %y, align 4
  %loaded_cond = load i32, ptr %y, align 4
  %cond_bool = icmp ne i32 %loaded_cond, 0
  br i1 %cond_bool, label %ternary.true, label %ternary.false

ternary.true:                                     ; preds = %entry
  br label %ternary.merge

ternary.false:                                    ; preds = %entry
  br label %ternary.merge

ternary.merge:                                    ; preds = %ternary.false, %ternary.true
  %ternarytmp = phi ptr [ %x, %ternary.true ], [ %y, %ternary.false ]
  %loaded_lhs = load i32, ptr %x, align 4
  %loaded_rhs = load i32, ptr %ternarytmp, align 4
  %gttmp = icmp sgt i32 %loaded_lhs, %loaded_rhs
  %select_result = select i1 %gttmp, ptr %x, ptr %y
  %selected_loaded = load i32, ptr %select_result, align 4
  store i32 %selected_loaded, ptr %max, align 4
  %loaded_int = load i32, ptr %max, align 4
  call void (ptr, ...) @printf(ptr @0, i32 %loaded_int)
  %ret = load i32, ptr %retval, align 4
  ret i32 %ret
}

declare void @printf(ptr, ...)

\end{lstlisting}



\begin{lstlisting}[
    caption={Сгенерированный LLVM IR код для тестовой программы №12},
    label={lst:while-ll},
    language={LLVM},
    basicstyle=\footnotesize\ttfamily,
    columns=fullflexible,
    keepspaces=true,
    breaklines=true
]

; ModuleID = 'generated'
source_filename = "generated"
target datalayout = "e-m:o-p270:32:32-p271:32:32-p272:64:64-i64:64-i128:128-n32:64-S128-Fn32"
target triple = "arm64-apple-darwin23.4.0"

@0 = private unnamed_addr constant [4 x i8] c"%d\0A\00", align 1
@1 = private unnamed_addr constant [2 x i8] c" \00", align 1
@2 = private unnamed_addr constant [4 x i8] c"%s\0A\00", align 1
@3 = private unnamed_addr constant [4 x i8] c"%d\0A\00", align 1

define i32 @main() {
entry:
  %retval = alloca i32, align 4
  %arr = alloca [5 x i32], align 4
  %arr_decay = getelementptr [5 x i32], ptr %arr, i32 0, i32 0
  %arr_elem = getelementptr i32, ptr %arr_decay, i32 0
  store i32 1, ptr %arr_elem, align 4
  %arr_elem1 = getelementptr i32, ptr %arr_decay, i32 1
  store i32 2, ptr %arr_elem1, align 4
  %arr_elem2 = getelementptr i32, ptr %arr_decay, i32 2
  store i32 3, ptr %arr_elem2, align 4
  %arr_elem3 = getelementptr i32, ptr %arr_decay, i32 3
  store i32 4, ptr %arr_elem3, align 4
  %arr_elem4 = getelementptr i32, ptr %arr_decay, i32 4
  store i32 5, ptr %arr_elem4, align 4
  %i = alloca i32, align 4
  store i32 0, ptr %i, align 4
  %loaded_from_var = load i32, ptr %i, align 4
  store i32 %loaded_from_var, ptr %i, align 4
  %j = alloca i32, align 4
  store i32 4, ptr %j, align 4
  %loaded_from_var5 = load i32, ptr %j, align 4
  store i32 %loaded_from_var5, ptr %j, align 4
  %tmpCounter = alloca i32, align 4
  store i32 0, ptr %tmpCounter, align 4
  %loaded_from_var6 = load i32, ptr %tmpCounter, align 4
  store i32 %loaded_from_var6, ptr %tmpCounter, align 4
  br label %while.cond

while.cond:                                       ; preds = %while.body, %entry
  %loaded_lhs = load i32, ptr %tmpCounter, align 4
  %lttmp = icmp slt i32 %loaded_lhs, 5
  %whilecond = icmp ne i1 %lttmp, false
  br i1 %whilecond, label %while.body, label %while.end

while.body:                                       ; preds = %while.cond
  %loaded_index = load i32, ptr %tmpCounter, align 4
  %arr_elem7 = getelementptr i32, ptr %arr_decay, i32 %loaded_index
  %arr_val = load i32, ptr %arr_elem7, align 4
  call void (ptr, ...) @printf(ptr @0, i32 %arr_val)
  %loadtmp = load i32, ptr %tmpCounter, align 4
  %inctmp = add i32 %loadtmp, 1
  store i32 %inctmp, ptr %tmpCounter, align 4
  br label %while.cond

while.end:                                        ; preds = %while.cond
  br label %while.cond8

while.cond8:                                      ; preds = %while.body9, %while.end
  %loaded_lhs11 = load i32, ptr %i, align 4
  %loaded_rhs = load i32, ptr %j, align 4
  %lttmp12 = icmp slt i32 %loaded_lhs11, %loaded_rhs
  %whilecond13 = icmp ne i1 %lttmp12, false
  br i1 %whilecond13, label %while.body9, label %while.end10

while.body9:                                      ; preds = %while.cond8
  %loaded_index14 = load i32, ptr %i, align 4
  %arr_elem15 = getelementptr i32, ptr %arr_decay, i32 %loaded_index14
  %arr_val16 = load i32, ptr %arr_elem15, align 4
  %tmp1 = alloca i32, align 4
  store i32 %arr_val16, ptr %tmp1, align 4
  %loaded_index17 = load i32, ptr %j, align 4
  %arr_elem18 = getelementptr i32, ptr %arr_decay, i32 %loaded_index17
  %arr_val19 = load i32, ptr %arr_elem18, align 4
  %tmp2 = alloca i32, align 4
  store i32 %arr_val19, ptr %tmp2, align 4
  %loaded_index20 = load i32, ptr %i, align 4
  %arr_elem21 = getelementptr i32, ptr %arr_decay, i32 %loaded_index20
  %arr_loaded_rhs = load i32, ptr %tmp2, align 4
  store i32 %arr_loaded_rhs, ptr %arr_elem21, align 4
  %loaded_index22 = load i32, ptr %j, align 4
  %arr_elem23 = getelementptr i32, ptr %arr_decay, i32 %loaded_index22
  %arr_loaded_rhs24 = load i32, ptr %tmp1, align 4
  store i32 %arr_loaded_rhs24, ptr %arr_elem23, align 4
  %loadtmp25 = load i32, ptr %i, align 4
  %inctmp26 = add i32 %loadtmp25, 1
  store i32 %inctmp26, ptr %i, align 4
  %loadtmp27 = load i32, ptr %j, align 4
  %dectmp = sub i32 %loadtmp27, 1
  store i32 %dectmp, ptr %j, align 4
  br label %while.cond8

while.end10:                                      ; preds = %while.cond8
  call void (ptr, ...) @printf(ptr @2, ptr @1)
  %tmpCounter28 = alloca i32, align 4
  store i32 0, ptr %tmpCounter28, align 4
  %loaded_from_var29 = load i32, ptr %tmpCounter28, align 4
  store i32 %loaded_from_var29, ptr %tmpCounter28, align 4
  br label %while.cond30

while.cond30:                                     ; preds = %while.body31, %while.end10
  %loaded_lhs33 = load i32, ptr %tmpCounter28, align 4
  %lttmp34 = icmp slt i32 %loaded_lhs33, 5
  %whilecond35 = icmp ne i1 %lttmp34, false
  br i1 %whilecond35, label %while.body31, label %while.end32

while.body31:                                     ; preds = %while.cond30
  %loaded_index36 = load i32, ptr %tmpCounter28, align 4
  %arr_elem37 = getelementptr i32, ptr %arr_decay, i32 %loaded_index36
  %arr_val38 = load i32, ptr %arr_elem37, align 4
  call void (ptr, ...) @printf(ptr @3, i32 %arr_val38)
  %loadtmp39 = load i32, ptr %tmpCounter28, align 4
  %inctmp40 = add i32 %loadtmp39, 1
  store i32 %inctmp40, ptr %tmpCounter28, align 4
  br label %while.cond30

while.end32:                                      ; preds = %while.cond30
  %ret = load i32, ptr %retval, align 4
  ret i32 %ret
}

declare void @printf(ptr, ...)

\end{lstlisting}